\documentclass[12pt]{article}

\usepackage{amsmath, amssymb, amsthm}
\usepackage{enumitem}
\usepackage{geometry}
\usepackage{fancyhdr} % for headers/footers
\usepackage{graphicx} % for pics
\usepackage{hyperref}

% --- Page Setup ---
\geometry{margin=1in}
\pagestyle{fancy}
\fancyhf{}
\rhead{MATH 421 -- 003, Fall 2025}
\lhead{Homework 3}
\cfoot{\thepage}
\renewcommand{\headrulewidth}{0.4pt}

% --- Title info ---
\title{Homework 3}
\author{Alejandro De La Torre \\ MATH 421 -- The Theory of Single Variable Calculus \\ Section 003 \\ Instructor: Grace Work}
\date{September 2025}

% --- Theorem environments ---
\newtheorem{theorem}{Theorem}
\newtheorem{lemma}{Lemma}
\newtheorem{proposition}{Proposition}
\newtheorem{corollary}{Corollary}

\begin{document}

\maketitle

\section*{Collaboration Statement}
I worked on this homework independently. No outside collaborators or resources were used, except for the textbook \textit{Elementary Analysis: The Theory of Calculus} by Ross and my own class notes. 

\section*{Problem 1}
Prove if $x$ and $y$ are real numbers, then:
\begin{enumerate}[label=\Alph*.]
  \item $|x-y| \le |x| + |y|$
  \item $|x| - |y| \le |x-y|$
\end{enumerate}

\textbf{Proof.}

\begin{enumerate}[label=\Alph*.]
\item We want to prove $|x-y| \le |x| + |y|$.

First write the difference as a sum:
\[
x-y = x + (-y).
\]
By the triangle inequality applied to the numbers $x$ and $-y$,
\[
|x-y|
= |x + (-y)|
\le |x| + |-y|
= |x| + |y|.
\]
Thus $|x-y| \le |x| + |y|$, as required.

\item We want to prove $|x| - |y| \le |x-y|$.

Start with the triangle inequality for the sum $(x-y)+y$:
\[
|(x-y)+y| \le |x-y| + |y|.
\]
The left-hand side simplifies to $|x|$, so
\[
|x| \le |x-y| + |y|.
\]
Subtract $|y|$ from both sides (this is valid because $|y|$ is a real number):
\[
|x| - |y| \le |x-y|.
\]
Hence $|x| - |y| \le |x-y|$, as desired.
\end{enumerate}
\hfill$\square$


\section*{Problem 2 (Ross Chapter 1, 3.6)}
\begin{enumerate}[label=\alph*.]

\item Prove $|a+b+c| \le |a|+|b|+|c|$ for all $a,b,c\in\mathbb{R}$. \\
\textit{Hint:} Apply the triangle inequality twice. Do not consider eight cases.

\begin{proof}
By associativity of addition,
\[
|a+b+c| \;=\; |(a+b)+c|.
\]
Applying the triangle inequality $|u+v|\le |u|+|v|$ with $u=a+b$ and $v=c$ gives
\[
|(a+b)+c| \;\le\; |a+b|+|c|.
\]
Apply the triangle inequality again to $|a+b|$:
\[
|a+b| \;\le\; |a|+|b|.
\]
Substituting this into the previous inequality yields
\[
|a+b+c| \;\le\; (|a|+|b|)+|c| \;=\; |a|+|b|+|c|,
\]
as required.
\end{proof}

\item Use induction to prove that for $n$ real numbers $a_1,\dots,a_n$,
\[
\bigl|a_1+a_2+\cdots+a_n\bigr| \;\le\; |a_1|+|a_2|+\cdots+|a_n|.
\]

\begin{proof}
Let $P(n)$ denote the statement
\[
\bigl|a_1+\cdots+a_n\bigr| \le |a_1|+\cdots+|a_n|
\quad\text{for all real }a_1,\dots,a_n.
\]

\textbf{Base case.}
For $n=1$, $|a_1| \le |a_1|$ is trivial.  
(If you prefer to start at $n=2$, then $|a_1+a_2|\le |a_1|+|a_2|$ is exactly the triangle inequality.)

\textbf{Inductive step.}
Assume $P(n)$ holds for some $n\ge 1$; i.e.,
\[
\bigl|a_1+\cdots+a_n\bigr| \le |a_1|+\cdots+|a_n|
\quad\text{for all }a_1,\dots,a_n\in\mathbb{R}.
\]
Let $a_1,\dots,a_n,a_{n+1}\in\mathbb{R}$ be arbitrary. Using associativity and the triangle inequality,
\begin{align*}
\bigl|a_1+\cdots+a_n+a_{n+1}\bigr|
&= \bigl|(a_1+\cdots+a_n)+a_{n+1}\bigr| \\
&\le \bigl|a_1+\cdots+a_n\bigr| + |a_{n+1}|.
\end{align*}
By the inductive hypothesis,
\[
\bigl|a_1+\cdots+a_n\bigr| \le |a_1|+\cdots+|a_n|.
\]
Combining the two inequalities gives
\[
\bigl|a_1+\cdots+a_n+a_{n+1}\bigr|
\le (|a_1|+\cdots+|a_n|)+|a_{n+1}|
= |a_1|+\cdots+|a_n|+|a_{n+1}|.
\]
Thus $P(n+1)$ holds.

By the principle of mathematical induction, $P(n)$ is true for every $n\in\mathbb{N}$.
\end{proof}

\end{enumerate}

\section*{Problem 3 (Ross Chapter 1, 3.8)}
Let $a,b\in\mathbb{R}$. Show: if $a\le b_1$ for every $b_1>b$, then $a\le b$.

\begin{proof}
We prove the contrapositive. Suppose $a>b$. Define
\[
b_1=\frac{a+b}{2}.
\]
Since $a>b$, we have $a-b>0$, hence $\dfrac{a-b}{2}>0$. Adding $b$ to both sides gives
\[
b+\frac{a-b}{2}=\frac{a+b}{2}=b_1 \;>\; b.
\]
Also, from $a-b>0$ we get $a-\dfrac{a-b}{2}=\dfrac{a+b}{2}=b_1<a$. Therefore $b_1>b$ and $a\not\le b_1$.

Thus, if $a>b$ there exists $b_1>b$ for which $a\not\le b_1$. This is the contrapositive of the desired implication, and hence the original statement holds: if $a\le b_1$ for every $b_1>b$, then $a\le b$.
\end{proof}

\section*{Problem 4}

\textbf{Definition.} An integer $n$ is divisible by an integer $d$ if $n = dk$ for some $k \in \mathbb{Z}$.

\medskip
\textbf{Theorem.} Suppose $n$ is an integer. If $n^2$ is divisible by $3$, then $n$ is divisible by $3$.

\begin{proof}
We proceed by proving the contrapositive: If $n$ is not divisible by $3$, then $n^2$ is not divisible by $3$.

If $n$ is not divisible by $3$, then by the Division Algorithm we can write
\[
n = 3k+1 \quad \text{or} \quad n = 3k+2
\]
for some integer $k$.

\textbf{Case 1.} If $n = 3k+1$, then
\[
n^2 = (3k+1)^2 = 9k^2 + 6k + 1 = 3(3k^2 + 2k) + 1.
\]
Thus $n^2 \equiv 1 \pmod{3}$, so $n^2$ is not divisible by $3$.

\textbf{Case 2.} If $n = 3k+2$, then
\[
n^2 = (3k+2)^2 = 9k^2 + 12k + 4 = 3(3k^2 + 4k + 1) + 1.
\]
Again, $n^2 \equiv 1 \pmod{3}$, so $n^2$ is not divisible by $3$.

\medskip
In both cases, if $n$ is not divisible by $3$, then $n^2$ is not divisible by $3$. By contrapositive, if $n^2$ is divisible by $3$, then $n$ must be divisible by $3$.
\end{proof}

\section*{Problem 5}
Prove that $\sqrt{3}$ is irrational.

\begin{proof}
Suppose, toward a contradiction, that $\sqrt{3}$ is rational. Then there exist
integers $p,q$ with $q\neq 0$ such that
\[
\sqrt{3}=\frac{p}{q}\qquad\text{and}\qquad \gcd(p,q)=1
\]
(i.e., the fraction is in lowest terms). Squaring both sides yields
\[
3=\frac{p^2}{q^2}\quad\Longrightarrow\quad p^2=3q^2.
\]
Hence $p^2$ is divisible by $3$. By the result of Problem~4 (if $n^2$ is divisible
by $3$, then $n$ is divisible by $3$), it follows that $3\mid p$. Thus there exists
$k\in\mathbb Z$ with $p=3k$. Substituting back, we obtain
\[
3q^2=p^2=(3k)^2=9k^2 \quad\Longrightarrow\quad q^2=3k^2,
\]
so $3\mid q$ as well.

Therefore $p$ and $q$ are both divisible by $3$, which contradicts $\gcd(p,q)=1$.
This contradiction shows our initial assumption was false. Hence $\sqrt{3}$ is
irrational.
\end{proof}

\section*{Problem 6}
Prove that $\sqrt{2}+\sqrt{3}$ and $\sqrt{2}-\sqrt{3}$ are both irrational numbers.

\begin{proof}
\textbf{(1) $\boldsymbol{\sqrt{2}+\sqrt{3}}$ is irrational.}
Suppose, for a contradiction, that $r=\sqrt{2}+\sqrt{3}\in\mathbb{Q}$.
Then $\sqrt{3}=r-\sqrt{2}$ and squaring gives
\[
3=(r-\sqrt{2})^2=r^2-2r\sqrt{2}+2.
\]
Rearranging yields
\[
-2r\sqrt{2}=1-r^2 \qquad\Longrightarrow\qquad
\sqrt{2}=\frac{r^2-1}{2r}.
\]
Since $r\in\mathbb{Q}$, the right-hand side is rational, which implies
$\sqrt{2}\in\mathbb{Q}$, a contradiction. Hence $\sqrt{2}+\sqrt{3}$ is irrational.

\medskip
\textbf{(2) $\boldsymbol{\sqrt{2}-\sqrt{3}}$ is irrational.}
Observe that
\[
(\sqrt{2}+\sqrt{3})(\sqrt{2}-\sqrt{3})=2-3=-1.
\]
If $\sqrt{2}-\sqrt{3}$ were rational, then
\(
\sqrt{2}+\sqrt{3}=-1/(\sqrt{2}-\sqrt{3})
\)
would be rational, contradicting part (1).
Therefore $\sqrt{2}-\sqrt{3}$ is irrational.
\end{proof}

\end{document}